\begin{tikzpicture}[
  >={Latex[length=3mm,width=2.2mm]},
  every node/.style={font=\small},
  uvec/.style={->,line width=1pt,black},
  mvec/.style={->,line width=1.3pt,vpurple},
  ivec/.style={->,line width=1pt,vpurple},
  bvec/.style={->,line width=1.3pt,bblue},
  dashline/.style={dashed,thin,gray!60!black}
]

% Spira (grigia, piccola, in basso a sinistra)
\draw[blobgray!80,thick] (-0.3,0) ellipse (0.6 and 0.15);

% Corrente i
\draw[ivec] (0.2,-0.15) arc (-90:180:0.3 and 0.15);
\node[vpurple,below] at (0,-0.3) {$i$};

% Vettore m (viola, verso l'alto)
\draw[mvec] (-0.3,0.2) -- ++(0,1.2) node[above left] {$\mathbf{m}$};

% Asse tratteggiato verticale
\draw[dashline] (-0.3,-0.5) -- (-0.3,1.5);

% u_r versore da spira verso P
\draw[uvec] (-0.2,0.3) -- ++(1.0,0.5) node[below right] {$\mathbf{u}_r$};

% Angolo theta
\draw[thin] (-0.3,0.8) arc (90:30:0.5);
\node at (-0.05,0.7) {$\theta$};

% Linea r
\draw[thin] (-0.3,0.3) -- (4.5,2.5);
\node[above] at (2.0,1.5) {$r$};

% Punto P in alto a destra
\coordinate (P) at (4.5,2.5);
\fill (P) circle (1.5pt);

% Parallelogramma giallo per decomposizione di B
\fill[ytangent,opacity=0.6] (P) -- ++(1.5,0) -- ++(-0.4,1.0) -- ++(-1.5,0) -- cycle;

% Vettori B al punto P (azzurri)
\draw[bvec] (P) -- ++(1.5,0) node[right] {$\mathbf{B}$};
\draw[bvec] (P) -- ++(-0.4,1.0) node[above left] {$\mathbf{B}_r$};
\draw[bvec] (P) -- ++(0.4,-0.8) node[below right] {$\mathbf{B}_\theta$};

\end{tikzpicture}
